\chapter{Introduction}
\label{chap:introduction}

\section{Problem Statement}
\label{sec:problem}

Automated Essay Scoring (AES) has entered a new era. Where earlier
systems counted surface features such as word length, sentence
complexity, and spelling, large language models (LLMs) now evaluate what previous
systems could not: the quality of a student's argument. Frontier models
are increasingly deployed to score high-stakes writing at scale,
promising consistent, rapid, and inexpensive assessment. Yet a model
that judges \emph{how well a student argues} must hold some standard of
what good argumentation looks like, and that standard is not neutral.
LLMs derive their evaluative expectations from training corpora that
overwhelmingly represent Western, Educated, Industrialised, Rich, and
Democratic (WEIRD) cultural contexts; empirical work has shown that
widely used models express cultural values closely aligned with
English-speaking and Protestant-European countries
\citep{tao2024cultural}.

Consider an illustration. A Vietnamese student opens an argumentative
essay with her grandfather's proverb, \emph{nuoc chay da mon} (water
wears down stone).
%% NOTE: proverb romanised deliberately -- mitthesis.cls under pdflatex
%% does not support Vietnamese diacritics without vntex/XeLaTeX.
She then develops her position the way many Asian
academic writing traditions teach: context first, evidence accumulated
gradually, the conclusion arrived at rather than announced. A trained
human examiner can recognise the coherence of this inductive structure.
An LLM evaluator calibrated by its training distribution to expect a
thesis statement in the opening paragraph may instead record ``lacks a
clear thesis'' and ``argument feels indirect''. It may also penalise the
hedging forms that signal politeness in her tradition as weak
reasoning. If so, she loses marks not because her English is poor, but
because the shape of her reasoning is culturally unfamiliar to the
model.

This thesis names this mechanism \emph{evaluative epistemic filtering}:
when an LLM evaluates writing, it applies an implicit cultural filter
derived from its training distribution, systematically discounting
rhetorical strategies that deviate from Anglo-American conventions
regardless of their argumentative merit. A pilot study in the preceding
phase of this research (Part~A) reported statistically significant
scoring disparities between Southeast Asian (SEA) and native-English
cohorts
of equivalent human-rated quality across two frontier models. The same
pilot reported that a culturally-diversified scoring panel reduced the
disparity. Those findings motivated, but do not yet establish, the
claims this thesis investigates at scale.

The problem this thesis addresses can therefore be stated plainly.
LLM-based essay scoring applies culturally partial standards,
invisibly and at scale, to students whose academic futures depend on
the result. The disparity is documented, but no validated framework
exists to mitigate it in evaluative settings. Nor does any calibrated
mechanism exist to detect the essays it affects. This thesis
quantifies the problem across models, builds and tests such a
framework, and tests its detection mechanism against human ground
truth.

\section{Research Questions and Hypotheses}
\label{sec:rqs}

Two research questions drive this thesis:

\begin{description}
  \item[RQ1.] To what extent does evaluative epistemic filtering
        manifest as measurable scoring disparities against Southeast
        Asian rhetorical patterns in LLM-driven essay scoring, and does
        it generalise across model families?
  \item[RQ2.] Can a culturally-aware multi-agent scoring architecture
        mitigate these disparities at inference time, without
        retraining or access to model weights, while reliably
        detecting the essays it cannot fairly score?
\end{description}

These questions are operationalised as four testable hypotheses:

\begin{itemize}
  \item \textbf{H1 (Generalisation).} Evaluative bias against Southeast
        Asian rhetorical patterns appears as a native$-$SEA cohort
        score gap whose 95\% confidence interval excludes zero on at
        least four LLMs, including an open-weight model.
  \item \textbf{H2 (Mitigation).} A five-agent culturally-calibrated
        panel with a meta-adjudicator reduces the cohort score gap
        more than the three-agent pilot configuration, without
        degrading agreement with the corpus's proficiency bands.
  \item \textbf{H3 (Detection).} A Bias Divergence Score threshold
        calibrated by receiver operating characteristic (ROC) analysis
        attains at least 90\% sensitivity at no more
        than a 5\% false-positive rate against multi-rater human
        disagreement on the 140-essay Global Rating Archives (GRA)
        validation set.
  \item \textbf{H4 (Trade-off boundary).} Consensus-weight allocation
        exhibits an identifiable boundary beyond which cohort-fairness
        gains degrade scoring accuracy against graded ground truth.
\end{itemize}

The research questions and hypotheses relate as follows. RQ1 is a
question about the \emph{existence and generality} of the problem; it
is operationalised by H1, which makes it falsifiable by fixing the
scope (four model families) and the test (statistical significance of
the cohort scoring gap). RQ2 is a question about the \emph{solution},
and it decomposes into three parts: whether the mitigation works (H2),
whether its failures can be detected (H3), and what the mitigation
costs (H4). Table~\ref{tab:rq-hypotheses} summarises the mapping and
where each hypothesis is tested.

\begin{table}[htbp]
\centering
\caption[Mapping of research questions to hypotheses]{How the research
questions decompose into hypotheses, and where each is tested.}
\label{tab:rq-hypotheses}
\small
\begin{tabular}{llll}
\hline
RQ & Hypothesis & Establishes & Tested in \\
\hline
RQ1 & H1 & Bias exists and generalises across models & Experiment E1 \\
RQ2 & H2 & The panel mitigates the bias & Experiments E1--E2 \\
RQ2 & H3 & Residual bias is detectable (BDS) & Experiment E3 \\
RQ2 & H4 & What mitigation costs in accuracy & Experiment E2 \\
\hline
\end{tabular}
\end{table}

Chapter~5 returns to each hypothesis explicitly. Because several
terms recur throughout, Table~\ref{tab:key-terms} defines them once
in plain language.

\begin{table}[htbp]
\centering
\caption[Key terms used in this thesis]{Key terms used in this
thesis, in plain language.}
\label{tab:key-terms}
\small
\begin{tabular}{p{3.3cm} p{8.9cm}}
\hline
Term & Meaning in this thesis \\
\hline
Evaluative epistemic filtering & The tendency of an AI grader to mark
down writing whose rhetorical form is unfamiliar from its training
data, regardless of the argument's quality. \\[2pt]
Lens & One AI evaluator given a documented cultural interpretation of
the same scoring rubric (Western, Southeast Asian, East Asian,
Mekong, or Neutral). \\[2pt]
Panel, consensus & Several lenses scoring the same essay
independently; their weighted average is the consensus score. \\[2pt]
Bias Divergence Score (BDS) & How much the lenses disagree on one
essay; high values are meant to flag the essay for a human. \\[2pt]
Cohort, cohort gap & A group of essays by writer background (Southeast
Asian learners or native-English writers); the gap is the difference
in mean score between the two cohorts. \\[2pt]
Band & The 1--9 integer score an evaluator assigns; the corpus's own
proficiency labels are CEFR bands (B1\_2, B2+). \\
\hline
\end{tabular}
\end{table}

\section{Theoretical Framing: Epistemic Injustice and Contrastive Rhetoric}
\label{sec:theory}

Two established bodies of theory ground this thesis. The first is
\citet{fricker2007epistemic}'s account of \emph{testimonial injustice}:
a speaker is wronged in their capacity as a knower when prejudice in
the hearer deflates the credibility of their testimony. Evaluative
epistemic filtering can be understood as testimonial injustice
automated and scaled: the student's argument is discounted not on its
merits but because its form does not match the evaluator's culturally
conditioned expectations. Fricker's framework also cautions against a
naive remedy: no evaluator, human or artificial, judges from nowhere.
This thesis therefore treats a ``culturally neutral'' evaluator as a
regulatory ideal rather than an achievable design goal, and pursues
\emph{plurality} of perspectives rather than pretended neutrality.

The second grounding is contrastive rhetoric. Since
\citet{kaplan1966cultural} first observed that rhetorical organisation
is culturally learned, research summarised by
\citet{connor2011intercultural} has documented systematic differences
between writing traditions. These include the distinction between
writer-responsible and reader-responsible coherence, and the validity of
inductive, context-first argument structures common in Asian academic
writing. This literature supplies the scholarly basis for the cultural
rubric interpretations at the centre of the proposed framework,
distinguishing them from ad-hoc cultural personas.

The mitigation strategy itself extends the formal multi-perspective
bias-reduction framework of \citet{dorleon2026uncovering}, who proved
that, under ideal conditions, generating text from multiple cultural
perspectives with bias filtering can eliminate stereotypical content
and perspective omission in \emph{generative} tasks. Companion
empirical work documented pervasive cultural bias in LLM treatments of
Southeast Asian content \citep{shujaa2026llms}. This
thesis asks whether those principles transfer from generation to
\emph{evaluation}. Independent support for the core mechanism comes
from the evaluation literature: panels of diverse LLM judges have been
shown to outperform and de-bias single large judges
\citep{verga2024replacing}, though with diversity defined by model
provider rather than by culture. The thesis, in short, tests whether
the jury effect holds when the axis of diversity is cultural.

Finally, the fairness--accuracy trade-off at the centre of H4 has a
direct lineage in fairness-aware machine learning.
\citet{dorleon2022feature} formulated feature selection as an explicit
trade-off between fairness and prediction performance. FAPFID
\citep{dorleon2023fapfid} then showed that fairness for unprivileged
groups can be improved without sacrificing overall performance,
through balanced representation. H4 asks whether an analogous frontier
exists in evaluative scoring, and where it lies.

\section{Significance}
\label{sec:significance}

The significance of this work is threefold. First, equity: as LLM-based
assessment becomes a gatekeeper to scholarships, university admission,
and migration pathways, a systematic penalty against non-Western
rhetorical traditions would constitute discrimination embedded in
infrastructure, invisible to its subjects and deniable by its
operators. Quantifying that penalty rigorously is a precondition for
governing it. Second, method: existing multi-agent AES research
improves accuracy without cultural awareness
\citep[e.g.,][]{su2025cafes}, while existing debiasing research targets
generative outputs rather than evaluative judgements; a validated
framework for \emph{evaluative} fairness would fill a documented gap at
their intersection. Third, deployability: because the proposed
framework operates entirely at inference time through commercial APIs,
any institution can apply it without retraining models. This practical
property distinguishes it from training-time approaches and
matches the constraints of real educational deployments, particularly
in the Southeast Asian context where this research is situated.

\section{Contributions}
\label{sec:contributions}


This thesis makes five contributions:

\begin{enumerate}
  \item \textbf{Conceptual:} it formalises evaluative epistemic
        filtering as a distinct bias mechanism in LLM-driven scoring,
        theoretically grounded in epistemic injustice and contrastive
        rhetoric.
  \item \textbf{Architectural:} it develops CAMPS (Culturally-Aware
        Multi-Perspective Scoring), a five-agent culturally-calibrated
        scoring panel with a meta-adjudicator, extending formal
        multi-perspective bias reduction from generation to evaluation.
  \item \textbf{Empirical:} it evaluates the framework on approximately
        200 topic-controlled essays from the ICNALE corpus
        \citep{ishikawa2023icnale} across four frontier and open-weight
        models, re-establishing the pilot evidence at scale on a corpus
        purpose-built for the comparison.
  \item \textbf{Methodological:} it introduces the Bias Divergence
        Score as a candidate trigger for targeted human review, and
        evaluates it against multi-rater human ground truth.
  \item \textbf{Practical:} it characterises the fairness--accuracy
        trade-off frontier, giving deploying institutions an explicit
        account of what fairness currently costs.
\end{enumerate}

\section{Thesis Outline}
\label{sec:outline}

The thesis follows the research questions in order. Chapter~2
reviews the literature that frames RQ1 and RQ2 and locates the gap
this thesis occupies. Chapter~3 develops the CAMPS framework and the
method: the data, the three experiments (E1 for H1 and part of H2,
E2 for H2 and H4, E3 for H3), the statistical protocol, and the alternatives
considered. Chapter~4 reports the experiments in that order and
discusses what the results mean, why they arose, how they compare
with prior work, and where they are uncertain. Chapter~5 answers each
hypothesis, states the contributions and limitations, and outlines
future work.
